% ID: SEL-D0040-OB-E03
% ESTADO: cotejada. Segunda pasada visual y borrador compilado revisados.
% Fuente: Selectividad/matematicas_aplicadas/Ua5MATEMÁTICAS APLICADAS examen.pdf
% SHA-256: ac80a164e049669cd01083ca5adaddfa90efbf556b050de2d99f294ab55670c9
% Página: 2; opción B; ejercicio original 3.

Una empresa dedicada a la producción de jamones ibéricos dispone de dos secaderos, A y B, con distintas condiciones ambientales y de almacenamiento. En el secadero B se curan la tercera parte de los jamones. El 25\% de los jamones curados en el secadero A son catalogados como Reserva, mientras que en el B este porcentaje asciende al 80\%.

Elegido un jamón al azar de uno de los secaderos, calcule la probabilidad de los siguientes sucesos:
\begin{enumerate}[label=\alph*)]
\item \textbf{(1.5 puntos)} El jamón no es de Reserva.
\item \textbf{(1 punto)} Si el jamón es de Reserva, que proceda del secadero A.
\end{enumerate}
